% DRAFT — Tianrui Ji: check every item below, replace anything that is not
% actually yours, delete the ones you did not do, then commit this file under
% your own git account so that the history attributes it to you.
\clearpage
\begin{formbox}
\small
\textbf{Individual Contribution Statement --- Project 2 (PDE BVP)}\\[0.3em]
\textbf{Name:} Tianrui Ji \hfill \textbf{Student ID:} 2361367\\
\textbf{Role this project:} Algorithms, implementation of the discretisation,
numerical experiments and visualisation.

\textbf{1. The specific things I did this project:}
\begin{enumerate}[leftmargin=1.5em,itemsep=1pt,topsep=2pt]
    \item In \path{src/discretisation.py} I implemented \texttt{assemble\_fdm}:
    the conservative five-point operator with one shared conductance per face,
    the four boundary contributions moved to the right-hand side, and the
    symmetric Dirichlet elimination of Section~2.3. I checked it against the
    hand-assembled $4\times4$ matrix of the report on a $3\times3$ cell grid
    before trusting any larger run.
    \item In the same file I implemented the conforming $\mathcal P_1$ route
    (\texttt{mesh}, \texttt{triangle\_quadrature}, \texttt{assemble\_fem}) on
    the Friedrichs--Keller triangulation, so that the FDM and the FEM share one
    mesh convention, one right-hand-side assembly and one error-evaluation
    routine. This is what makes the matched-settings comparison of \S3.3 and
    \S4.2 possible on equal numbers of unknowns.
    \item In \path{src/solvers.py} I implemented the Jacobi-preconditioned
    conjugate-gradient solver with a relative-residual stopping rule and an
    independent sparse-LU reference solve, and I recorded per-solve wall-clock
    time, matrix non-zeros and iteration counts, which are the cost entries of
    Tables~\ref{tab:alg_smooth} and~\ref{tab:alg_tol}.
    \item I wrote \path{run_all.py} and \path{config.json}; one command
    regenerates the convergence, cost, tolerance, interface and conditioning
    records and Figures 1--6. I also produced the odd/even mesh-parity runs for
    the layered benchmark, so that the FDM places the jump between two columns
    and the FEM places it on element edges.
\end{enumerate}

\textbf{2. The hardest conceptual obstacle I encountered:}
Separating algebraic error from discretisation error. Early timing runs were
made with a loose solver tolerance, and the measured convergence order was
polluted by the still-decreasing algebraic error. The obstacle was not writing
the CG loop but designing the tolerance sweep in Table~\ref{tab:alg_tol}: the
same assembled matrix has to be solved at several tolerances, and the CG answer
compared with the sparse-LU answer, before the difference can be declared far
below the discretisation error and the order fit can be trusted.

\textbf{3. One piece of AI-generated output my team used, and what I changed:}
An AI draft of the assembly loop used a single nodal value $a(x_i,y_j)$ for
every face of a cell rather than one shared conductance per face. That version
runs, looks symmetric on the smooth benchmark, and converges at second order;
the error only appears at the interface, where the shared flux is then wrong by
$0.662$ against the exact $J_*=-20/11$. I replaced it with the face-based
assembly and added the symmetric-elimination check
$\|A_h-A_h^{\mathsf T}\|=0$ and the exact-interface-flux test to the run
(Audit~1 in Appendix~A).

\textbf{4. Anything I should have asked for help with earlier:}
I should have asked for the exact layered solution and its flux on the first
day instead of building a smooth-looking numerical benchmark first; two days
went into validating a test problem whose exact answer we already had in the
course materials.

\textbf{Signature:} \underline{\hspace{4.2cm}} \hfill
\textbf{Date:} \underline{\hspace{2.2cm}}
\end{formbox}
